\documentclass[11pt,a4paper]{article}
\usepackage[utf8]{inputenc}
\usepackage[T1]{fontenc}
\usepackage[french]{babel}
\usepackage{amsmath,amssymb,amsthm}
\usepackage{bm}
\usepackage{mathtools}
\usepackage{geometry}
\usepackage{hyperref}
\geometry{margin=2.5cm}

\title{Formalisation mathematique d'un modele hybride\\HyConEx--HyperLogic}
\author{Projet de recherche}
\date{\today}

\begin{document}
\maketitle

\begin{abstract}
Ce document formalise mathematiquement une architecture hybride combinant
l'approche contrefactuelle de HyConEx et l'approche de generation de regles de type HyperLogic.
L'objectif est d'unifier:
(i) une tete de prediction interpretable par regles,
(ii) une tete de generation de contrefactuels,
et (iii) un hyperreseau generateur de poids.
Nous definissons les notations, les equations de passage avant, la fonction de cout multi-objectifs,
et les criteres d'evaluation.
\end{abstract}

\section{Notations}
Soit un jeu de donnees supervises:
\[
\mathcal{D}=\{(\bm{x}_i,y_i)\}_{i=1}^N,\quad
\bm{x}_i\in\mathbb{R}^{d},\quad
y_i\in\{1,\dots,C\}.
\]

On note:
\begin{itemize}
    \item $B(\cdot)$ l'operateur de binarisation/discretisation;
    \item $\bm{x}^{(b)}=B(\bm{x})\in\{0,1\}^{D_b}$ la representation binaire;
    \item $\bm{t}\in\{1,\dots,C\}$ une classe cible contrefactuelle, $\bm{t}\neq y$;
    \item $\mathcal{H}_\phi$ l'hyperreseau de parametres $\phi$;
    \item $W_m$ les poids du main network (tete regles), $W_{cf}$ ceux de la tete CF.
\end{itemize}

\section{Structure hybride HyConEx--HyperLogic}
\subsection{Hyperreseau}
Dans l'hybride, les poids des deux tetes sont produits par un hyperreseau:
\[
(W_m,W_{cf})=\mathcal{H}_\phi(\bm{s}),
\]
ou $\bm{s}$ est un vecteur de contexte (par exemple statistiques de batch de $\bm{x}^{(b)}$:
moyenne, variance, quantiles, etc.).

Contrairement a une fusion de type $\alpha W_{hyper}+(1-\alpha)W_{base}$,
on considere ici la forme \emph{full generated}:
\[
W_m=W_m^{hyper},\qquad W_{cf}=W_{cf}^{hyper}.
\]

\subsection{Main network de regles}
Le main network est une fonction differentiable:
\[
\bm{z}=f_m(\bm{x}^{(b)};W_m),\qquad
\hat{\bm{p}}=\text{softmax}(\bm{z})\in[0,1]^C.
\]
La prediction est:
\[
\hat{y}=\arg\max_{c\in\{1,\dots,C\}} \hat{p}_c.
\]

Interpretation regles:
les neurones caches representent des clauses candidates et la couche de sortie agrège ces clauses.
Les litteraux sont derives des poids de $W_m$ et des variables binaires de $\bm{x}^{(b)}$.

\paragraph{Formalisation explicite de la tete regles}
Une ecriture compacte de type DR-like est:
\begin{align}
\bm{h} &= \sigma\!\left((\bm{x}^{(b)})^\top W_r + \bm{b}_r\right),\quad \bm{h}\in\mathbb{R}^{K},\\
\bm{z} &= \bm{h}^\top W_o + \bm{b}_o,\quad \bm{z}\in\mathbb{R}^{C},
\end{align}
avec:
\[
W_m=\{W_r,\bm{b}_r,W_o,\bm{b}_o\},
\]
ou $K$ est le nombre de neurones-regles, $W_r\in\mathbb{R}^{D_b\times K}$, $W_o\in\mathbb{R}^{K\times C}$.
La fonction $\sigma$ est typiquement une approximation differentiable d'une logique de conjonction
(ex. sigmoide temperaturee).

Pour un neurone-regle $k$, les conditions ON/OFF sont lues depuis la colonne $W_r[:,k]$:
un poids positif favorise l'activation du bit, un poids negatif favorise sa desactivation.

\paragraph{Lien direct avec la formulation HyperLogic (niveau neurone-regle)}
Dans HyperLogic/DR-Net, la sortie logique d'une regle $k$ est idealement definie par une indicatrice:
\begin{equation}
o_k
=\mathbf{1}\!\left\{
\sum_{d=1}^{D_b} w_{d,k}x^{(b)}_d-\sum_{d=1}^{D_b}|w_{d,k}|=0
\right\},
\label{eq:ok_indicator}
\end{equation}
ou $x^{(b)}_d$ est la composante binaire et $w_{d,k}$ le poids associe.
L'equation \eqref{eq:ok_indicator} encode l'idee ``tous les litteraux requis par la regle sont satisfaits''.

Comme \eqref{eq:ok_indicator} est non differentiable, on utilise une approximation lisse:
\begin{equation}
o_k \approx h(u_k),
\quad
u_k=\sum_{d=1}^{D_b} w_{d,k}x^{(b)}_d-\sum_{d=1}^{D_b}|w_{d,k}|,
\quad
h(u)=\exp\!\left(-\frac{u^2}{\tau}\right),\ \tau>0,
\label{eq:ok_smooth}
\end{equation}
ou $\tau$ controle la nettete logique (plus $\tau$ est petit, plus l'approximation se rapproche
d'une condition dure).

La couche de sortie de type OR/additif s'ecrit:
\begin{equation}
z_c=\sum_{k=1}^{K} u_{k,c}\,o_k + b_c,\qquad c=1,\dots,C.
\label{eq:or_layer}
\end{equation}
Ainsi, chaque classe combine les activations de regles via les coefficients $u_{k,c}$.

\paragraph{Version hyperreseau multi-experts}
Dans l'esprit HyperLogic, l'hyperreseau peut produire plusieurs jeux de poids
$\{(W_m^{(m)},W_{cf}^{(m)})\}_{m=1}^{M}$.
Pour la tete regles seule, on obtient:
\begin{equation}
f_M(\bm{x}^{(b)})=
\frac{1}{M}\sum_{m=1}^{M} f_m(\bm{x}^{(b)};W_m^{(m)}),
\label{eq:mc_experts}
\end{equation}
qui est une moyenne d'experts de regles.
Dans la pratique du projet, on peut prendre $M=1$ (poids uniques par batch) ou $M>1$
pour augmenter la diversite/robustesse des regles.

\subsection{Tete contrefactuelle}
La tete CF genere un exemple alternatif:
\[
\bm{x}_{cf}^{(b)}=g_{cf}\!\left(\bm{x}^{(b)},\text{onehot}(\bm{t});W_{cf}\right).
\]
On repasse ensuite par la tete de regles:
\[
\bm{z}_{cf}=f_m(\bm{x}_{cf}^{(b)};W_m),\qquad
\hat{\bm{p}}_{cf}=\text{softmax}(\bm{z}_{cf}),\qquad
\hat{y}_{cf}=\arg\max_c \hat{p}_{cf,c}.
\]

\paragraph{Formalisation explicite de la tete contrefactuelle}
On concatene entree binaire et cible:
\[
\bm{u}=[\bm{x}^{(b)};\text{onehot}(\bm{t})]\in\mathbb{R}^{D_b+C}.
\]
Une parametrisation a deux couches:
\begin{align}
\bm{q} &= \rho\!\left(\bm{u}^\top W_{cf,1}+\bm{b}_{cf,1}\right),\\
\bm{\pi} &= \text{sigmoid}\!\left(\bm{q}^\top W_{cf,2}+\bm{b}_{cf,2}\right),\quad \bm{\pi}\in(0,1)^{D_b},
\end{align}
avec:
\[
W_{cf}=\{W_{cf,1},\bm{b}_{cf,1},W_{cf,2},\bm{b}_{cf,2}\}.
\]
Le binaire de sortie peut etre echantillonne/seuille, ou traite en \emph{straight-through}:
\[
\bm{x}_{cf}^{(b)} \approx \text{ST}\!\left(\bm{\pi}\right).
\]
Cette tete apprend donc une politique de \emph{flip} de bits conditionnee par la classe cible.

\section{Formulation de la perte}
On combine quatre termes:
\[
\mathcal{L}
=\lambda_{cls}\mathcal{L}_{cls}
+\lambda_{cf}\mathcal{L}_{cf}
+\lambda_{flip}\mathcal{L}_{flip}
+\lambda_{rule}\mathcal{R}_{rule}.
\]

\subsection{Perte de classification}
\[
\mathcal{L}_{cls}
=\frac{1}{|\mathcal{B}|}\sum_{i\in\mathcal{B}}
\text{CE}\!\left(\hat{\bm{p}}_i,y_i\right).
\]

\subsection{Perte de validite contrefactuelle}
\[
\mathcal{L}_{cf}
=\frac{1}{|\mathcal{B}|}\sum_{i\in\mathcal{B}}
\text{CE}\!\left(\hat{\bm{p}}_{cf,i},t_i\right).
\]

\subsection{Parcimonie des modifications}
En espace binaire:
\[
\mathcal{L}_{flip}
=\frac{1}{|\mathcal{B}|}\sum_{i\in\mathcal{B}}
\left\|\bm{x}^{(b)}_{cf,i}-\bm{x}^{(b)}_i\right\|_1.
\]
Ce terme force des contrefactuels proches et actionnables.

\subsection{Regularisation de regles}
Une regularisation de sparsite sur les poids de la partie regles:
\[
\mathcal{R}_{rule}=\|W_m\|_1
\quad\text{(ou sous-ensemble structure de }W_m\text{)}.
\]
Objectif: reduire la complexite des regles extraites.

\section{Optimisation}
L'apprentissage consiste a resoudre:
\[
\phi^\star=\arg\min_{\phi}\;\mathbb{E}_{(\bm{x},y)\sim\mathcal{D}}[\mathcal{L}].
\]

Par descente de gradient stochastique (Adam/AdamW):
\[
\phi\leftarrow\phi-\eta\nabla_\phi\mathcal{L}.
\]

Comme $W_m,W_{cf}$ dependent de $\phi$, le gradient retro-propage via la composition:
\[
\phi\rightarrow(W_m,W_{cf})\rightarrow
\{f_m,g_{cf}\}\rightarrow\mathcal{L}.
\]

\section{Inference et explication}
\subsection{Prediction}
Pour une entree $\bm{x}$:
\[
\bm{x}^{(b)}=B(\bm{x}),\quad
(W_m,W_{cf})=\mathcal{H}_\phi(\bm{s}),\quad
\hat{y}=\arg\max f_m(\bm{x}^{(b)};W_m).
\]

\subsection{Contrefactuel}
Pour une cible $t\neq \hat{y}$:
\[
\bm{x}_{cf}^{(b)}=g_{cf}\!\left(\bm{x}^{(b)},\text{onehot}(t);W_{cf}\right),\quad
\hat{y}_{cf}=\arg\max f_m(\bm{x}_{cf}^{(b)};W_m).
\]
La validite est verifiee si $\hat{y}_{cf}=t$.

\subsection{Projection continue optionnelle}
Si necessaire pour l'usage metier:
\[
\bm{x}_{cf}=B^{-1}(\bm{x}_{cf}^{(b)}).
\]
Cette etape ne modifie pas la logique de decision du modele, elle sert a l'interpretation humaine.

\section{Mesures d'evaluation mathematiques}
\subsection{Classification}
\[
\text{Accuracy}
=\frac{1}{N}\sum_{i=1}^{N}\mathbf{1}(\hat{y}_i=y_i),
\]
avec AUROC OvR et matrice de confusion en complement.

\subsection{Qualite contrefactuelle}
\[
\text{Validity}
=\frac{1}{N}\sum_{i=1}^{N}\mathbf{1}(\hat{y}_{cf,i}=t_i),
\]
\[
\text{FlipMean}
=\frac{1}{N}\sum_{i=1}^{N}\left\|\bm{x}^{(b)}_{cf,i}-\bm{x}^{(b)}_i\right\|_1.
\]
Apres projection continue:
\[
\text{L1Cont}
=\frac{1}{N}\sum_{i=1}^{N}\|\bm{x}_{cf,i}-\bm{x}_i\|_1.
\]

\section{Lien conceptuel avec HyConEx et HyperLogic}
L'hybride peut etre lu comme:
\begin{itemize}
    \item \textbf{Heritage HyConEx}: objectif conjoint classification+contrefactuels;
    \item \textbf{Heritage HyperLogic}: hyperreseau producteur de poids pour un moteur de regles differentiable;
    \item \textbf{Synthese}: un unique schema d'optimisation qui preserve explicabilite et recourse.
\end{itemize}

\section{Conclusion}
Cette formalisation fournit une base mathematique unifiee pour
construire, entrainer et evaluer un modele hybride HyConEx--HyperLogic.
Le cadre est extensible:
regularisation de diversite des regles, contraintes de causalite sur les CF,
ou variantes multiclasses plus structurees.

\end{document}
